\begingroup\small\singlespacing
\setlength{\tabcolsep}{3pt}
\renewcommand{\arraystretch}{1.03}
\begin{longtable}{@{}>{\raggedright\arraybackslash}p{1.9cm}>{\raggedright\arraybackslash}p{3.05cm}>{\raggedright\arraybackslash}p{2.15cm}>{\raggedright\arraybackslash}p{2.9cm}>{\raggedright\arraybackslash}p{3.0cm}@{}}
\caption{Síntesis del cumplimiento y alcance de los objetivos. Los títulos abreviados remiten a los objetivos aprobados en la Sección~\ref{sec:objetivos}. Elaboración propia.}\label{tab:sintesis-objetivos}\\
\toprule Objetivo & Evidencia obtenida & Ubicación & Alcance efectivo & Limitaciones de la evidencia \\\midrule\endfirsthead
\multicolumn{5}{l}{\tablename~\thetable\ (continuación)}\\\toprule
Objetivo & Evidencia obtenida & Ubicación & Alcance efectivo & Limitaciones de la evidencia \\\midrule\endhead
\midrule\multicolumn{5}{r}{Continúa en la página siguiente}\\\endfoot
\bottomrule\endlastfoot
General: comparación reproducible & Predicciones, métricas y costos de seis configuraciones; modelos y registros conservados. & Secciones\newline\ref{sec:comparacion-clasificacion}--\ref{sec:compromiso}. & Clasificación de ModelNet40 en dos resoluciones, sobre $2\,468$ objetos de prueba. & Un equipo y una semilla; los alcances de algunos tiempos difieren. \\ \addlinespace
1. Enfoque clásico basado en octree & Representaciones construidas; equivalencia y persistencia verificadas en la muestra; clasificación clásica evaluada. & Secciones~\ref{sec:resultados-representaciones} y~\ref{sec:resultados-hce}; Tablas~\ref{tab:geometria-resultados} y~\ref{tab:construccion-costos}. & Inventario de $12\,311$ octrees por resolución; diez objetos y veinte registros en la validación controlada. & La verificación geométrica de la muestra no equivale a una auditoría exhaustiva del inventario. \\ \addlinespace
2. HCE, SVM y Bosque Aleatorio & Vectores de 17 y 18 componentes; contratos auditados; modelos seleccionados y persistidos. & Sección~\ref{sec:resultados-hce}; Tabla~\ref{tab:clasificacion-final}. & Partición de $8\,858$ objetos de ajuste y 985 de validación; prueba oficial separada. & Asociación fila--objeto respaldada por registros y revisión del programa, sin repetir localmente su ejecución. \\ \addlinespace
3. Net5 basado en octrees & Dos modelos entrenados; historiales y mejores estados conservados; evaluación final. & Sección~\ref{sec:resultados-net5}; Figura~\ref{fig:curvas-finales}; Tabla~\ref{tab:clasificacion-final}. & Mismas particiones documentadas; resoluciones $32^3$ y $64^3$. & Un entrenamiento por configuración; sin estudio de ablación ni variabilidad entre semillas. \\ \addlinespace
4. Protocolo reproducible y trazabilidad & Manifiestos, relaciones HCE, huellas y versiones; seis secuencias de predicción preservadas en la reevaluación. & Secciones~\ref{sec:resultados-hce} y~\ref{sec:comparacion-costo}; Tabla~\ref{tab:reevaluacion-final}. & Repetición de la evaluación de los modelos guardados; 103 pruebas aprobadas según el registro de ejecución. & No constituye reentrenamiento independiente; varían tiempos y memoria entre ejecuciones. \\ \addlinespace
5. Desempeño, costo y escalabilidad & Exactitud, confusiones, intervalos, contrastes, entrenamiento, latencia, memoria y tamaño. & Secciones\newline\ref{sec:comparacion-clasificacion}--\ref{sec:compromiso}; Tablas~\ref{tab:latencia-final}--\ref{tab:reevaluacion-final}; Anexo~\ref{anexo:confusiones}. & Prueba: $2\,468$ objetos; latencia: 400; memoria: 100; comparación entre dos resoluciones. & Extracción HCE estimada; memoria de tres categorías; tiempos con alcances distintos y sin medición de energía. \\
\end{longtable}\endgroup
